\chapter{Cotangent Spaces}
\label{ch:vii05-cotangent-spaces}

Tangent vectors differentiate functions. Cotangent vectors are the linear measurements that can be applied to tangent vectors:
\[
\boxed{T_p^*M=(T_pM)^*,\qquad df_p(v)=v(f).}
\]
They are the natural home of differentials, constraints, and the covariant side of coordinate calculus.

\section{Covectors and the cotangent space}
\begin{definition}[Cotangent space]\label{def:vii05-cotangent-space}
The cotangent space at \(p\) is the dual vector space
\[
T_p^*M=\operatorname{Hom}_{\mathbb R}(T_pM,\mathbb R).
\]
Its elements are called cotangent vectors or covectors.
\end{definition}

\begin{proposition}[Dimension]\label{prop:vii05-dimension}
If \(M\) is \(n\)-dimensional, then \(\dim T_p^*M=n\).
\end{proposition}
\begin{proof}
The dual of an \(n\)-dimensional vector space has dimension \(n\).
\end{proof}

\section{Coordinate coframes}
Let \(x=(x^1,\dots,x^n)\) be local coordinates near \(p\).

\begin{definition}[Coordinate covectors]\label{def:vii05-coordinate-covectors}
Define \(dx^i|_p\in T_p^*M\) by
\[
dx^i|_p(v)=v(x^i).
\]
\end{definition}

\begin{theorem}[Dual coordinate basis]\label{thm:vii05-dual-basis}
The family
\[
dx^1|_p,\dots,dx^n|_p
\]
is dual to
\[
\partial_{x^1}|_p,\dots,\partial_{x^n}|_p,
\]
so
\[
dx^i|_p\!\left(\left.\frac{\partial}{\partial x^j}\right|_p\right)=\delta^i_j.
\]
Every \(\alpha\in T_p^*M\) has the unique expansion
\[
\alpha=\sum_i \alpha\!\left(\left.\frac{\partial}{\partial x^i}\right|_p\right)dx^i|_p.
\]
\end{theorem}

\section{Differentials of functions}
\begin{definition}[Differential of a function]\label{def:vii05-df}
For \(f\in C^\infty(M)\), define
\[
df_p\in T_p^*M,\qquad df_p(v)=v(f).
\]
\end{definition}

\begin{theorem}[Coordinate formula for \(df\)]\label{thm:vii05-df-coordinate}
In local coordinates,
\[
df_p=\sum_{i=1}^n
\frac{\partial(f\circ x^{-1})}{\partial x^i}(x(p))\,dx^i|_p.
\]
\end{theorem}
\begin{proof}
Evaluate both sides on the coordinate basis vectors of \(T_pM\).
\end{proof}

\begin{proposition}[Algebra rules]\label{prop:vii05-df-rules}
For smooth functions \(f,g\) and \(a,b\in\mathbb R\),
\[
d(af+bg)=a\,df+b\,dg,\qquad
d(fg)=f\,dg+g\,df
\]
pointwise.
\end{proposition}
\begin{proof}
These identities are the linearity and Leibniz rule for tangent derivations.
\end{proof}

\section{Pullback of covectors}
\begin{definition}[Pullback at a point]\label{def:vii05-pullback-point}
For smooth \(F:M\to N\) and \(p\in M\), define
\[
F_p^*:T_{F(p)}^*N\to T_p^*M,\qquad
(F_p^*\alpha)(v)=\alpha(dF_pv).
\]
\end{definition}

\begin{theorem}[Pullback functoriality]\label{thm:vii05-pullback-functorial}
For \(F:M\to N\), \(G:N\to P\),
\[
(G\circ F)_p^*=F_p^*\circ G_{F(p)}^*,
\qquad
(\operatorname{id}_M)_p^*=\operatorname{id}_{T_p^*M}.
\]
\end{theorem}
\begin{proof}
This is dual to the tangent-space chain rule; evaluate on an arbitrary tangent vector.
\end{proof}

\begin{theorem}[Pulling back differentials]\label{thm:vii05-pullback-df}
For \(g\in C^\infty(N)\),
\[
F_p^*(dg_{F(p)})=d(g\circ F)_p.
\]
\end{theorem}
\begin{proof}
For \(v\in T_pM\), both sides evaluate to \(v(g\circ F)\).
\end{proof}

\section{Coordinate transformation law}
\begin{theorem}[Coframe transformation]\label{thm:vii05-coframe-transform}
For two coordinate systems \(x\) and \(y\),
\[
dy^a|_p=\sum_i\frac{\partial y^a}{\partial x^i}(p)\,dx^i|_p.
\]
Equivalently, covector components transform by the inverse transpose relative to the transformation law for tangent-vector components.
\end{theorem}
\begin{proof}
Since \(dy^a=d(y^a)\), apply the coordinate formula for the differential of the function \(y^a\).
\end{proof}

\begin{remark}[Covariant versus contravariant]\label{rem:vii05-covariant}
The basis covectors and their component columns transform oppositely so that the scalar pairing \(\alpha(v)\) remains invariant.
\end{remark}

\section{Annihilators}
\begin{definition}[Annihilator]\label{def:vii05-annihilator}
For a linear subspace \(W\subset T_pM\), its annihilator is
\[
W^\circ=\{\alpha\in T_p^*M:\alpha(w)=0\text{ for every }w\in W\}.
\]
\end{definition}

\begin{theorem}[Dimension of the annihilator]\label{thm:vii05-annihilator-dimension}
If \(\dim T_pM=n\) and \(\dim W=k\), then
\[
\dim W^\circ=n-k.
\]
\end{theorem}
\begin{proof}
Extend a basis of \(W\) to one of \(T_pM\) and inspect the dual basis.
\end{proof}

\begin{proposition}[Kernel annihilator]\label{prop:vii05-kernel-annihilator}
For a linear map \(A:V\to W\),
\[
\operatorname{im}(A^*)=(\ker A)^\circ.
\]
In finite dimensions, \(A\) is surjective iff \(A^*\) is injective, and \(A\) is injective iff \(A^*\) is surjective.
\end{proposition}

\section{Differentials and level sets}
If \(f:M\to\mathbb R\), the covector \(df_p\) measures first-order change of \(f\). When \(df_p\neq0\), its kernel is a hyperplane:
\[
\ker df_p=\{v\in T_pM:v(f)=0\}.
\]
VII/06 identifies this kernel with the tangent space to a regular level set.

\begin{definition}[Critical and regular points]\label{def:vii05-critical}
A point \(p\) is critical for \(f\) when \(df_p=0\); otherwise it is regular.
\end{definition}

\section{Differential versus gradient}
\begin{warning}[No gradient without extra structure]\label{warning:vii05-gradient}
The differential \(df_p\) is defined on every smooth manifold. A gradient vector \(\nabla f(p)\) requires a Riemannian metric \(g\) to identify covectors with vectors through
\[
g_p(\nabla f(p),v)=df_p(v).
\]
Without a metric, \(df_p\) and a tangent vector are different kinds of objects.
\end{warning}

\section{Examples}
\begin{example}[Euclidean differential]\label{ex:vii05-euclidean}
For \(f(x,y)=x^2y+\sin y\),
\[
df=(2xy)\,dx+(x^2+\cos y)\,dy.
\]
\end{example}

\begin{example}[Constraint covector]\label{ex:vii05-sphere-constraint}
For \(F(x)=\|x\|^2\) on \(\mathbb R^n\),
\[
dF_x=2\sum_i x^i\,dx^i.
\]
At \(x\neq0\), \(\ker dF_x\) is the Euclidean hyperplane orthogonal to \(x\); VII/06 recognizes it as \(T_xS^{n-1}\) on a sphere level.
\end{example}

\section{Exercises}
\begin{exercise}\label{exr:vii05-01}
Prove the cotangent space has the same dimension as the tangent space.
\end{exercise}
\begin{hint}
Use finite-dimensional duality.
\end{hint}
\begin{solution}
A vector space and its dual have the same finite dimension.
\end{solution}

\begin{exercise}\label{exr:vii05-02}
Verify \(dx^i|_p(\partial_{x^j}|_p)=\delta^i_j\).
\end{exercise}
\begin{hint}
Use the definition of coordinate derivations.
\end{hint}
\begin{solution}
\(\partial x^i/\partial x^j=\delta^i_j\).
\end{solution}

\begin{exercise}\label{exr:vii05-03}
Expand an arbitrary covector in the coordinate coframe.
\end{exercise}
\begin{hint}
Use the dual-basis theorem.
\end{hint}
\begin{solution}
\(\alpha=\sum_i\alpha(\partial_{x^i}|_p)dx^i|_p\).
\end{solution}

\begin{exercise}\label{exr:vii05-04}
Show \(df_p(v)=v(f)\) agrees with VII/04's differential \(df_p:T_pM\to T_{f(p)}\mathbb R\).
\end{exercise}
\begin{hint}
Identify \(T_{f(p)}\mathbb R\) with \(\mathbb R\).
\end{hint}
\begin{solution}
The two definitions give the same scalar on every tangent vector.
\end{solution}

\begin{exercise}\label{exr:vii05-05}
Compute \(df\) for \(f(x,y,z)=xe^{yz}\).
\end{exercise}
\begin{hint}
Take partial derivatives.
\end{hint}
\begin{solution}
\(df=e^{yz}dx+xze^{yz}dy+xye^{yz}dz\).
\end{solution}

\begin{exercise}\label{exr:vii05-06}
Prove linearity of the differential in the function.
\end{exercise}
\begin{hint}
Apply a tangent derivation.
\end{hint}
\begin{solution}
\(d(af+bg)_p(v)=v(af+bg)=a\,v(f)+b\,v(g)\).
\end{solution}

\begin{exercise}\label{exr:vii05-07}
Prove the product rule for differentials.
\end{exercise}
\begin{hint}
Use the derivation Leibniz rule.
\end{hint}
\begin{solution}
\(d(fg)=f\,dg+g\,df\) pointwise.
\end{solution}

\begin{exercise}\label{exr:vii05-08}
Derive the coordinate transformation law for coframes.
\end{exercise}
\begin{hint}
Apply the formula for \(d(y^a)\).
\end{hint}
\begin{solution}
\(dy^a=\sum_i(\partial y^a/\partial x^i)dx^i\).
\end{solution}

\begin{exercise}\label{exr:vii05-09}
Show the scalar pairing between a covector and tangent vector is chart independent.
\end{exercise}
\begin{hint}
Combine the dual transformation laws.
\end{hint}
\begin{solution}
The transition Jacobian and its inverse cancel in the pairing.
\end{solution}

\begin{exercise}\label{exr:vii05-10}
Prove \(F_p^*\) is linear.
\end{exercise}
\begin{hint}
Use linearity of a covector.
\end{hint}
\begin{solution}
\(F_p^*(a\alpha+b\beta)(v)=a\alpha(dF_pv)+b\beta(dF_pv)\).
\end{solution}

\begin{exercise}\label{exr:vii05-11}
Prove pullback reverses composition.
\end{exercise}
\begin{hint}
Evaluate on a tangent vector.
\end{hint}
\begin{solution}
\((G\circ F)^*=F^*G^*\) follows from \(d(G\circ F)=dG\,dF\).
\end{solution}

\begin{exercise}\label{exr:vii05-12}
Prove \(F^*(dg)=d(g\circ F)\) at each point.
\end{exercise}
\begin{hint}
Evaluate on a tangent vector.
\end{hint}
\begin{solution}
Both sides give the directional derivative of the composite.
\end{solution}

\begin{exercise}\label{exr:vii05-13}
For \(F(x,y)=(x+y,x-y)\) and \(\alpha=du+2\,dv\), compute \(F^*\alpha\).
\end{exercise}
\begin{hint}
Compute \(d(u\circ F)\) and \(d(v\circ F)\).
\end{hint}
\begin{solution}
\(F^*\alpha=(dx+dy)+2(dx-dy)=3dx-dy\).
\end{solution}

\begin{exercise}\label{exr:vii05-14}
Find the annihilator of \(W=\operatorname{span}(e_1,\dots,e_k)\subset\mathbb R^n\).
\end{exercise}
\begin{hint}
Use the dual standard basis.
\end{hint}
\begin{solution}
\(W^\circ=\operatorname{span}(dx^{k+1},\dots,dx^n)\).
\end{solution}

\begin{exercise}\label{exr:vii05-15}
Prove the annihilator dimension theorem.
\end{exercise}
\begin{hint}
Extend a basis and dualize.
\end{hint}
\begin{solution}
If \(W\) has basis \(e_1,\dots,e_k\), its annihilator is spanned by the remaining dual basis vectors.
\end{solution}

\begin{exercise}\label{exr:vii05-16}
Show \(\operatorname{im}(A^*)\subseteq(\ker A)^\circ\).
\end{exercise}
\begin{hint}
Evaluate \(A^*\alpha\) on \(v\in\ker A\).
\end{hint}
\begin{solution}
\((A^*\alpha)(v)=\alpha(Av)=0\).
\end{solution}

\begin{exercise}\label{exr:vii05-17}
Prove equality \(\operatorname{im}(A^*)=(\ker A)^\circ\) in finite dimensions.
\end{exercise}
\begin{hint}
Compare dimensions after proving inclusion.
\end{hint}
\begin{solution}
Both dimensions equal \(\operatorname{rank}A\).
\end{solution}

\begin{exercise}\label{exr:vii05-18}
Show a surjective \(dF_p\) gives an injective pullback \(F_p^*\).
\end{exercise}
\begin{hint}
Use linear duality.
\end{hint}
\begin{solution}
The dual of a surjective finite-dimensional linear map is injective.
\end{solution}

\begin{exercise}\label{exr:vii05-19}
Show an injective \(dF_p\) gives a surjective pullback \(F_p^*\).
\end{exercise}
\begin{hint}
Use linear duality.
\end{hint}
\begin{solution}
The dual of an injective finite-dimensional linear map is surjective.
\end{solution}

\begin{exercise}\label{exr:vii05-20}
For \(f(x,y)=x^2+y^2\), determine critical points.
\end{exercise}
\begin{hint}
Solve \(df=0\).
\end{hint}
\begin{solution}
\(df=2x\,dx+2y\,dy\), so only \((0,0)\) is critical.
\end{solution}

\begin{exercise}\label{exr:vii05-21}
At \(p\neq0\), compute \(\ker d(\|x\|^2)_p\).
\end{exercise}
\begin{hint}
Write the differential.
\end{hint}
\begin{solution}
It is \(\{v:p\cdot v=0\}\).
\end{solution}

\begin{exercise}\label{exr:vii05-22}
Explain why a gradient cannot be defined canonically on a smooth manifold from its smooth structure alone.
\end{exercise}
\begin{hint}
A covector-to-vector identification is missing.
\end{hint}
\begin{solution}
The smooth structure gives the dual space but no preferred isomorphism between a vector space and its dual.
\end{solution}

\begin{exercise}\label{exr:vii05-23}
Given a Riemannian metric \(g\), show the equation \(g(\nabla f,\cdot)=df\) determines \(\nabla f\) uniquely.
\end{exercise}
\begin{hint}
Use nondegeneracy of the inner product.
\end{hint}
\begin{solution}
The musical map from tangent vectors to covectors is an isomorphism.
\end{solution}

\begin{exercise}\label{exr:vii05-24}
State the bridge from annihilators to regular level sets.
\end{exercise}
\begin{hint}
Think of \(\ker dF_p\).
\end{hint}
\begin{solution}
For a regular level set \(S=F^{-1}(c)\), VII/06 proves \(T_pS=\ker dF_p\); equivalently the conormal space is generated by the pulled-back coordinate differentials of the constraints.
\end{solution}


\section{Solved problem dossiers}
\begin{problem}[Legacy cotangent problem: dual basis]\label{prob:vii05-01}
Construct the coordinate coframe and prove it is dual to the coordinate tangent basis.
\end{problem}
\begin{solution}
Define \(dx^i(v)=v(x^i)\); evaluation on \(\partial_{x^j}\) gives \(\delta^i_j\), so dual-basis theory completes the proof.
\end{solution}

\begin{problem}[Legacy differential problem: \(df\)]\label{prob:vii05-02}
Show that the differential of a smooth real-valued function is intrinsically a cotangent vector.
\end{problem}
\begin{solution}
Define \(df_p(v)=v(f)\). This is linear in \(v\), hence belongs to \(T_p^*M\), and its coordinate formula follows by evaluating on the coordinate basis.
\end{solution}

\begin{problem}[Legacy covector problem: transformation law]\label{prob:vii05-03}
Derive the coordinate change law for covectors.
\end{problem}
\begin{solution}
Use \(dy^a=d(y^a)=\sum_i(\partial y^a/\partial x^i)dx^i\), then dualize to obtain the component transformation law.
\end{solution}

\begin{problem}[Pullback construction]\label{prob:vii05-04}
Construct the pullback of a covector under a smooth map and prove it is coordinate free.
\end{problem}
\begin{solution}
Set \(F_p^*\alpha=\alpha\circ dF_p\). Both ingredients are intrinsic, so no chart enters the definition.
\end{solution}

\begin{problem}[Pullback chain rule]\label{prob:vii05-05}
Prove pullback reverses composition.
\end{problem}
\begin{solution}
Dualizing \(d(G\circ F)=dG\circ dF\) gives \((G\circ F)^*=F^*\circ G^*\).
\end{solution}

\begin{problem}[Differential of a composite]\label{prob:vii05-06}
Prove \(d(g\circ F)=F^*(dg)\).
\end{problem}
\begin{solution}
Evaluate both covectors on an arbitrary tangent vector and use the definition of the differential.
\end{solution}

\begin{problem}[Annihilator dimension]\label{prob:vii05-07}
For \(W\subset T_pM\), prove \(\dim W+\dim W^\circ=\dim M\).
\end{problem}
\begin{solution}
Extend a basis of the subspace to a tangent basis and inspect the dual basis.
\end{solution}

\begin{problem}[Kernel-image duality]\label{prob:vii05-08}
Prove \(\operatorname{im}(A^*)=(\ker A)^\circ\).
\end{problem}
\begin{solution}
Inclusion is immediate; rank-nullity gives equal dimensions in finite dimensions.
\end{solution}

\begin{problem}[Critical-point criterion]\label{prob:vii05-09}
Show \(p\) is critical for \(f\) iff every tangent direction has zero first-order change.
\end{problem}
\begin{solution}
This is exactly \(df_p=0\iff df_p(v)=v(f)=0\) for all \(v\).
\end{solution}

\begin{problem}[Constraint hyperplane]\label{prob:vii05-10}
For a regular scalar constraint \(f=c\), explain why \(\ker df_p\) has codimension one.
\end{problem}
\begin{solution}
A nonzero linear functional on an \(n\)-dimensional vector space has an \((n-1)\)-dimensional kernel.
\end{solution}

\begin{problem}[Differential versus gradient]\label{prob:vii05-11}
Separate the notions of differential and gradient and state the extra datum needed for the latter.
\end{problem}
\begin{solution}
The differential is a covector defined from smooth structure alone. A Riemannian metric supplies the isomorphism \(v\mapsto g(v,\cdot)\) needed to turn \(df\) into \(\nabla f\).
\end{solution}

\begin{problem}[Diffeomorphism and cotangent spaces]\label{prob:vii05-12}
Show a diffeomorphism induces cotangent-space isomorphisms in the reverse direction.
\end{problem}
\begin{solution}
If \(F\) is a diffeomorphism, \(dF_p\) is an isomorphism; its dual \(F_p^*\) is therefore an isomorphism with inverse \((F^{-1})_{F(p)}^*\).
\end{solution}

\section{Challenge problems}
\begin{problem}[Cotangent space from first-order function germs]\label{prob:vii05-ch01}
Show \(\mathfrak m_p/\mathfrak m_p^2\cong T_p^*M\).
\end{problem}
\begin{solution}
Coordinate germs \(x^i-x^i(p)\) form a basis of the quotient and map to \(dx^i|_p\). This is dual to the tangent-space description from VII/04.
\end{solution}

\begin{problem}[Double annihilator]\label{prob:vii05-ch02}
For \(W\subset T_pM\), prove \((W^\circ)^\circ=W\) under the canonical identification with the double dual.
\end{problem}
\begin{solution}
The inclusion is immediate. Dimensions agree: the double annihilator has dimension \(n-(n-\dim W)=\dim W\).
\end{solution}

\begin{problem}[Conormal space of several constraints]\label{prob:vii05-ch03}
If \(F=(f^1,\dots,f^k):M\to\mathbb R^k\) has rank \(k\) at \(p\), show \(df^1_p,\dots,df^k_p\) are linearly independent.
\end{problem}
\begin{solution}
They are the pullbacks under \(dF_p^*\) of the standard dual basis. Surjectivity of \(dF_p\) makes the dual map injective.
\end{solution}

\begin{problem}[Invariant pairing]\label{prob:vii05-ch04}
Prove directly that tangent and cotangent coordinate transformation laws leave the pairing invariant.
\end{problem}
\begin{solution}
Write the transformations with a Jacobian and inverse Jacobian; their matrix product is the identity.
\end{solution}

\begin{problem}[Metric identification is additional structure]\label{prob:vii05-ch05}
Show two different inner products on one vector space can send the same covector to different gradient vectors.
\end{problem}
\begin{solution}
On \(\mathbb R^2\), take \(df=dx\). The Euclidean metric gives gradient \((1,0)\), while the metric with matrix \(\operatorname{diag}(2,1)\) gives \((1/2,0)\).
\end{solution}

\section*{Bridge to VII/06}
Tangent kernels and cotangent annihilators now describe constraints intrinsically. VII/06 uses the inverse and implicit function principles to turn regular constraints into embedded submanifolds and proves \(T_pF^{-1}(c)=\ker dF_p\), while product manifolds give the canonical splitting \(T_{(p,q)}(M\times N)\cong T_pM\oplus T_qN\).
